\lecture{9.2}{Multidimensional Exposures}
\label{lec:9-2}

\subsection{Radiation Oncology Example}

Some exposures are naturally high-dimensional. In radiation oncology, the
treatment is not simply ``radiation'' versus ``no radiation.'' A patient receives
a spatial dose distribution across tissues. For head and neck cancer, dose to
salivary glands may affect xerostomia, food intake, and weight loss. Minor
changes in where dose is delivered can matter clinically even if the total dose
is similar.

Let $A$ denote a vector of radiation-dose summaries and let $Y$ denote a
secondary outcome such as severe weight loss. A causal question may be:
what would the mean outcome be if every patient received dose map $a$ rather
than dose map $a'$?
\[
  \E\{\Ya{a}\}-\E\{\Ya{a'}\}.
\]
The difficulty is that $a$ and $a'$ may be high-dimensional objects rather than
simple scalar doses.

\subsection{3D Dose Maps}

A raw radiation plan can be represented as a three-dimensional voxel map:
\[
  A=(A_1,\ldots,A_p),
\]
where $A_j$ is the dose delivered to voxel $j$. The dimension $p$ may be very
large. The causal response surface
\[
  a\mapsto \E\{\Ya{a}\}
\]
is therefore a function over a high-dimensional exposure space.

\begin{figure}[H]
\centering
\begin{tikzpicture}[
  cell/.style={draw, minimum size=.42cm},
  arr/.style={-{Latex[length=2.2mm]}, thick}
]
\foreach \x in {0,...,4} {
  \foreach \y in {0,...,3} {
    \pgfmathtruncatemacro{\shade}{20+8*\x+10*\y}
    \node[cell, fill=blue!\shade] at (.42*\x,.42*\y) {};
  }
}
\node[right=2.0cm of current bounding box.east] (vec) {$A=(A_1,\ldots,A_p)$};
\draw[arr] (1.1,.8) -- (vec);
\end{tikzpicture}
\caption{A spatial exposure such as radiation dose can be encoded as a high-dimensional vector.}
\label{fig:ch9-dose-map}
\end{figure}

\subsection{Dose-Volume Summaries}

Clinical analyses often replace the full map with lower-dimensional summaries,
such as the dose received by each percentile of a gland's volume. For example,
$A_k$ might be the dose at the $k$th percentile, $k=1,\ldots,100$. This
summarizes localization but still leaves a high-dimensional ordered vector.

Such summaries are useful, but they are not automatically causal. If two dose
maps have the same summary but different biological effects, the summary loses
causally relevant information. If the summary is associated with confounders
such as tumor location or disease severity, associational prediction may not
recover causal effects.

\subsection{High-Dimensional Text Exposures}

In natural language processing applications, the exposure may be a document,
message, or clinical note. The intervention might alter wording, tone, content,
or information ordering. Text can be encoded as word counts, embeddings, topics,
or transformer features, but each representation changes the meaning of the
intervention.

For causal inference, one must specify what it means to intervene on text. A
change in a word embedding is not automatically a meaningful intervention unless
it corresponds to an interpretable change in text generation.

\subsection{Neuroimaging Exposures}

Neuroimaging studies may treat patterns of neural activity, connectivity, or
regional activation as exposures. These exposures are structured, correlated,
and high-dimensional. A causal claim might concern the effect of activity in a
network on behavior, but manipulating one region while holding the rest fixed
may be biologically impossible.

\subsection{Environmental Mixtures}

Environmental epidemiology often studies mixtures: particulate matter, ozone,
heavy metals, temperature, humidity, and other exposures occurring jointly.
The causal effect of one pollutant may depend on the levels of others, and the
positivity condition may fail because some mixtures are never observed.

\subsection{Why Multidimensional Treatments Complicate Causal Interpretation}

Multidimensional exposures create several linked problems:
\begin{itemize}
  \item \textbf{Estimand complexity}: the response surface
  $\E\{\Ya{a}\}$ is indexed by a high-dimensional $a$.
  \item \textbf{Intervention meaning}: not every vector $a$ corresponds to a
  feasible intervention.
  \item \textbf{Positivity}: the observed data may cover only a small region of
  the exposure space.
  \item \textbf{Confounding}: high-dimensional exposures are often assigned by
  complex clinical, environmental, or social processes.
  \item \textbf{Interpretability}: even if a model predicts well, the estimated
  causal contrast may not map to an actionable intervention.
\end{itemize}

\begin{keyidea}
High-dimensional exposure analysis begins with defining a meaningful
intervention. Dimension reduction and machine learning are useful only after the
causal estimand is scientifically interpretable.
\end{keyidea}
